\section{农业上市公司明细表}
\label{app:companies}

本附录列出申万农林牧渔行业（801010）全部 \num{104} 家 A 股上市公司的画像数据。
第一张表为\hl{上市公司画像总表}：财务指标统一采用 \textbf{2026 年半年报}（合并口径、未年化），
资产负债率为 2026 年 6 月末数值，市值时点为 2026 年 9 月 24 日；
“再融资/重组并购”两列为公司近三年（2023 年 9 月---2026 年 9 月）相应类别公告的数量，
用于衡量公司在融资与并购上的活跃程度。
第二张表按申万三级行业汇总各子行业的公司数量与整体经营水平，口径与第一张表一致；
第三张表\hl{债务结构与滚动偿付压力}（\ref{tab:company-financing}）把 2026 年 6 月末的
货币资金、短期债务（短期借款 $+$ 一年内到期的非流动负债）、有息负债
（短期债务 $+$ 长期借款 $+$ 应付债券）与\emph{现金短债比}列在一起，按现金短债比升序排列，
并在末列标出第 \ref{sec:financing-need} 节的分层归属；“静态缺口”为
$\max(0,\ \text{短期债务} - \text{货币资金})$，即假设不续贷、不动用其他资产时的滚动偿付缺口。
需要强调的是，短期债务在实务中大多通过银行续贷滚动，账面货币资金也可能包含受限资金与
募集资金专户，因此该表用于公司之间的横向比较，不能相加为全行业的“待融资总额”。

\begingroup\scriptsize\setlength{\tabcolsep}{1.5pt}
\begin{longtable}{@{}llp{1.9cm}rrrrrr@{}}
\caption{申万农林牧渔行业 A 股上市公司画像总表（按流通市值降序；市值时点 2026-09-24，财务为 2026 年半年报（未年化），资产负债率为 2026 年 6 月末）}\label{tab:company-master}\\
\toprule
代码 & 公司简称 & 申万三级行业 & 市值（亿元） & 营收（亿元） & 同比（\%） & 净利（亿元） & ROE（\%） & 负债率（\%） \\
\midrule
\endfirsthead
\multicolumn{9}{l}{\small（续）}\\
\toprule
代码 & 公司简称 & 申万三级行业 & 市值（亿元） & 营收（亿元） & 同比（\%） & 净利（亿元） & ROE（\%） & 负债率（\%） \\
\midrule
\endhead
\bottomrule
\endlastfoot
002714 & 牧原股份 & 生猪养殖 & 1,327.9 & 594.1 & -22.3 & -60.78 & -7.5 & 54.2 \\
300498 & 温氏股份 & 生猪养殖 & 862.7 & 467.6 & -6.2 & -43.66 & -11.1 & 58.9 \\
002311 & 海大集团 & 水产饲料 & 680.4 & 642.4 & 9.2 & 16.11 & 6.5 & 54.6 \\
600737 & 中粮糖业 & 其他农产品加工 & 315.5 & 92.8 & -21.2 & 6.08 & 5.2 & 39.1 \\
000876 & 新希望 & 畜禽饲料 & 311.4 & 555.9 & 7.7 & -17.02 & -9.1 & 70.7 \\
601118 & 海南橡胶 & 其他种植业 & 277.3 & 151.2 & -25.4 & -1.30 & -1.3 & 70.7 \\
002157 & 正邦科技 & 生猪养殖 & 227.4 & 77.2 & 12.9 & -7.36 & -7.0 & 53.5 \\
600598 & 北大荒 & 粮食种植 & 222.0 & 21.3 & -29.4 & -5.37 & -7.1 & 39.5 \\
002299 & 圣农发展 & 肉鸡养殖 & 191.3 & 107.8 & 21.7 & 5.02 & 4.4 & 51.9 \\
605198 & 安德利 & 果蔬加工 & 172.9 & 7.2 & -24.4 & 1.23 & 4.3 & 7.8 \\
000592 & 平潭发展 & 林业Ⅲ & 168.1 & 4.8 & -35.1 & -0.69 & -4.0 & 39.3 \\
605296 & 神农集团 & 生猪养殖 & 156.9 & 24.9 & -11.0 & -8.09 & -17.9 & 40.7 \\
300999 & 金龙鱼 & 粮油加工 & 133.2 & 1,228.5 & 6.2 & 22.94 & 2.4 & 56.2 \\
300761 & 立华股份 & 肉鸡养殖 & 132.1 & 91.3 & 9.3 & -0.72 & -0.8 & 43.8 \\
601952 & 苏垦农发 & 粮食种植 & 125.5 & 43.3 & -5.6 & 1.44 & 2.0 & 46.5 \\
600201 & 生物股份 & 动物保健Ⅲ & 125.4 & 7.6 & 22.1 & 1.04 & 1.9 & 16.5 \\
000998 & 隆平高科 & 种子 & 114.1 & 17.1 & -21.1 & -2.73 & -4.1 & 59.1 \\
002385 & 大北农 & 畜禽饲料 & 106.4 & 138.1 & 1.9 & -6.77 & -9.3 & 69.9 \\
600127 & 金健米业 & 粮油加工 & 105.8 & 19.7 & 25.2 & -0.10 & -1.4 & 62.2 \\
002100 & 天康生物 & 畜禽饲料 & 103.5 & 83.3 & -5.9 & -4.41 & -6.6 & 55.6 \\
000048 & 京基智农 & 生猪养殖 & 98.8 & 21.3 & -10.7 & -4.89 & -12.9 & 63.6 \\
000930 & 中粮科技 & 其他农产品加工 & 97.6 & 81.9 & -7.0 & -1.69 & -1.6 & 40.7 \\
002891 & 中宠股份 & 宠物食品 & 89.6 & 32.8 & 34.9 & 1.27 & 4.0 & 44.0 \\
002041 & 登海种业 & 种子 & 87.5 & 4.7 & 27.0 & 0.50 & 1.3 & 15.8 \\
600108 & 亚盛集团 & 其他种植业 & 80.4 & 14.7 & -2.6 & 0.25 & 0.6 & 58.9 \\
002458 & 益生股份 & 肉鸡养殖 & 78.2 & 17.0 & 28.4 & 3.08 & 6.9 & 35.1 \\
600251 & 冠农股份 & 果蔬加工 & 77.9 & 13.9 & 0.3 & 5.37 & 13.3 & 30.8 \\
002215 & 诺普信 & 其他种植业 & 77.7 & 39.6 & 7.8 & 7.65 & 16.8 & 60.5 \\
301498 & 乖宝宠物 & 宠物食品 & 75.2 & 35.4 & 10.0 & 1.91 & 4.1 & 31.4 \\
600313 & 农发种业 & 种子 & 73.6 & 38.7 & 32.8 & 0.53 & 2.2 & 37.2 \\
603477 & 巨星农牧 & 生猪养殖 & 70.8 & 35.6 & -4.3 & -8.38 & -30.3 & 69.4 \\
600962 & 国投中鲁 & 果蔬加工 & 68.6 & 8.8 & -12.4 & 0.02 & 0.2 & 52.0 \\
600195 & 中牧股份 & 动物保健Ⅲ & 67.8 & 31.7 & 13.5 & 0.50 & 0.9 & 30.2 \\
000735 & 罗牛山 & 生猪养殖 & 64.9 & 12.4 & 14.5 & -3.10 & -7.4 & 52.3 \\
600975 & 新五丰 & 生猪养殖 & 63.5 & 28.3 & -23.2 & -7.36 & -41.0 & 84.2 \\
603336 & 宏辉果蔬 & 其他种植业 & 62.7 & 4.2 & -11.1 & 0.05 & 0.4 & 36.3 \\
603668 & 天马科技 & 水产饲料 & 61.6 & 31.3 & 5.3 & 0.12 & 0.6 & 72.1 \\
603363 & 傲农生物 & 畜禽饲料 & 60.1 & 45.2 & 14.3 & -1.90 & -6.2 & 62.1 \\
688526 & 科前生物 & 动物保健Ⅲ & 56.8 & 3.7 & -23.3 & 1.16 & 2.8 & 10.7 \\
300119 & 瑞普生物 & 动物保健Ⅲ & 56.4 & 17.2 & 0.8 & 2.04 & 4.3 & 34.4 \\
603609 & 禾丰股份 & 畜禽饲料 & 56.2 & 189.6 & 8.9 & -0.08 & -0.1 & 57.7 \\
300189 & 神农种业 & 种子 & 54.7 & 0.6 & -35.3 & -0.13 & -1.9 & 39.8 \\
002567 & 唐人神 & 畜禽饲料 & 51.3 & 108.8 & -12.7 & -10.46 & -24.9 & 71.1 \\
001313 & 粤海饲料 & 水产饲料 & 51.1 & 37.8 & 41.7 & 0.48 & 1.9 & 52.1 \\
600354 & 敦煌种业 & 种子 & 50.3 & 8.7 & 20.9 & 1.45 & 18.7 & 37.5 \\
300087 & ST荃银 & 种子 & 49.9 & 11.8 & -17.9 & -0.55 & -3.3 & 64.8 \\
002321 & 华英农业 & 其他养殖 & 48.2 & 23.7 & 14.5 & -0.29 & -3.2 & 66.2 \\
002556 & 辉隆股份 & 农业综合Ⅱ & 47.4 & 85.5 & 3.3 & 1.13 & 3.1 & 67.0 \\
002481 & 双塔食品 & 其他农产品加工 & 46.5 & 11.2 & 7.3 & 0.33 & 1.4 & 44.8 \\
600371 & 万向德农 & 种子 & 46.4 & 1.0 & -15.1 & 0.14 & 2.6 & 14.5 \\
001338 & 永顺泰 & 其他农产品加工 & 45.7 & 19.1 & -13.9 & 0.51 & 1.4 & 21.1 \\
300138 & 晨光生物 & 其他农产品加工 & 44.6 & 31.6 & -13.6 & 1.79 & 5.2 & 64.1 \\
600191 & 华资实业 & 其他农产品加工 & 43.9 & 3.5 & 58.0 & 0.36 & 2.1 & 23.8 \\
002746 & 仙坛股份 & 肉鸡养殖 & 43.9 & 29.9 & 17.9 & 1.80 & 3.9 & 43.0 \\
300871 & 回盛生物 & 动物保健Ⅲ & 42.2 & 8.6 & 5.2 & 2.06 & 9.1 & 31.8 \\
000505 & 京粮控股 & 粮油加工 & 42.1 & 31.1 & -26.0 & 0.08 & 0.3 & 43.5 \\
301116 & 益客食品 & 肉鸡养殖 & 40.9 & 99.0 & 13.2 & 0.34 & 2.2 & 68.9 \\
002124 & 天邦食品 & 生猪养殖 & 40.8 & 36.6 & -22.1 & -15.43 & -96.8 & 91.2 \\
002772 & 众兴菌业 & 食用菌 & 39.5 & 10.9 & 19.6 & 1.91 & 5.3 & 54.0 \\
000713 & 国投丰乐 & 种子 & 39.0 & 11.9 & 3.9 & -0.31 & -1.0 & 32.1 \\
002688 & 金河生物 & 动物保健Ⅲ & 38.7 & 14.5 & 4.1 & 1.31 & 5.7 & 56.1 \\
000798 & 中水渔业 & 海洋捕捞 & 36.5 & 20.8 & 18.9 & 1.02 & 20.9 & 71.6 \\
603566 & 普莱柯 & 动物保健Ⅲ & 35.3 & 4.9 & -12.2 & 0.77 & 3.0 & 15.2 \\
600467 & 好当家 & 水产养殖 & 34.8 & 8.0 & 13.0 & 0.30 & 0.9 & 49.1 \\
688098 & 申联生物 & 动物保健Ⅲ & 34.4 & 1.1 & -8.1 & -0.19 & -1.4 & 9.8 \\
300021 & 大禹节水 & 农业综合Ⅱ & 33.4 & 15.9 & 24.3 & -0.01 & -0.0 & 71.7 \\
002086 & 东方海洋 & 其他农产品加工 & 32.3 & 1.6 & 0.9 & -0.48 & -3.8 & 19.6 \\
300094 & 国联水产 & 水产养殖 & 32.1 & 14.5 & -12.5 & 0.10 & 10.0 & 95.2 \\
600540 & 新赛股份 & 其他种植业 & 31.6 & 25.7 & 77.2 & 1.25 & 30.8 & 82.7 \\
603151 & 邦基科技 & 畜禽饲料 & 31.4 & 26.8 & 11.8 & -0.00 & -0.0 & 63.4 \\
002286 & 保龄宝 & 其他农产品加工 & 31.3 & 15.0 & 6.9 & 0.82 & 3.6 & 23.1 \\
001201 & 东瑞股份 & 生猪养殖 & 31.2 & 10.4 & -5.1 & -4.25 & -14.3 & 52.9 \\
002548 & 金新农 & 畜禽饲料 & 30.7 & 23.2 & -2.4 & -3.88 & -37.4 & 81.5 \\
000972 & 中基健康 & 果蔬加工 & 30.5 & 2.0 & -19.6 & -0.66 & — & 104.1 \\
600265 & *ST景谷 & 林业Ⅲ & 30.4 & 0.6 & -54.6 & -0.00 & -0.0 & 26.7 \\
600359 & 新农开发 & 其他种植业 & 29.1 & 2.0 & -31.1 & 0.19 & 2.8 & 55.6 \\
000019 & 深粮控股 & 粮油加工 & 28.2 & 23.9 & 0.1 & 1.12 & 2.2 & 38.1 \\
300511 & 雪榕生物 & 食用菌 & 27.6 & 11.6 & 46.2 & 0.18 & 1.1 & 56.5 \\
002679 & 福建金森 & 林业Ⅲ & 27.3 & 0.5 & -9.5 & -0.14 & -1.8 & 61.9 \\
600257 & 大湖股份 & 水产养殖 & 26.4 & 4.3 & 2.0 & 0.17 & 2.1 & 46.4 \\
002069 & 獐子岛 & 水产养殖 & 26.1 & 7.3 & -5.3 & 0.13 & 29.5 & 97.0 \\
300175 & 朗源股份 & 果蔬加工 & 25.2 & 1.7 & 67.7 & -0.13 & -2.7 & 8.8 \\
300313 & 天山生物 & 其他养殖 & 24.4 & 2.2 & 427.7 & 0.04 & 12.1 & 82.0 \\
002868 & 绿康生化 & 动物保健Ⅲ & 24.1 & 2.5 & -15.7 & 0.07 & 10.9 & 90.1 \\
300673 & 佩蒂股份 & 宠物食品 & 23.8 & 7.3 & 1.0 & 0.24 & 1.4 & 35.3 \\
002234 & 民和股份 & 肉鸡养殖 & 23.4 & 10.3 & 4.2 & -0.95 & -5.5 & 53.9 \\
002852 & 道道全 & 粮油加工 & 23.1 & 28.6 & 2.6 & 0.17 & 0.8 & 49.1 \\
688156 & 路德科技 & 畜禽饲料 & 23.0 & 2.1 & 42.9 & 0.09 & 1.2 & 55.9 \\
603182 & 嘉华股份 & 粮油加工 & 22.2 & 8.6 & 25.7 & 0.57 & 5.0 & 18.8 \\
002696 & 百洋股份 & 水产饲料 & 22.2 & 20.5 & 21.1 & -0.83 & -6.2 & 66.4 \\
600097 & 开创国际 & 海洋捕捞 & 22.2 & 11.8 & -5.4 & -0.30 & -1.3 & 34.4 \\
600965 & *ST福成 & 其他养殖 & 22.0 & 5.7 & 6.9 & 0.18 & 0.8 & 16.7 \\
003030 & 祖名股份 & 其他农产品加工 & 21.9 & 10.1 & 10.0 & -0.24 & -2.5 & 57.9 \\
300970 & 华绿生物 & 食用菌 & 21.8 & 6.9 & 39.2 & 0.46 & 2.8 & 35.4 \\
001366 & 播恩集团 & 畜禽饲料 & 21.0 & 6.7 & 9.0 & 0.10 & 1.3 & 38.1 \\
000911 & *ST广糖 & 其他农产品加工 & 20.5 & 15.8 & 17.4 & -0.16 & 0.0 & 100.2 \\
300967 & 晓鸣股份 & 肉鸡养殖 & 20.2 & 5.7 & -24.5 & 0.02 & 0.2 & 52.3 \\
000702 & 正虹科技 & 畜禽饲料 & 19.7 & 5.7 & 14.2 & -0.21 & -6.2 & 54.2 \\
600883 & 博闻科技 & 食用菌 & 19.1 & 0.3 & -8.8 & 0.15 & 1.5 & 6.6 \\
000663 & 永安林业 & 林业Ⅲ & 18.6 & 1.5 & 14.7 & -0.19 & -2.1 & 36.6 \\
603231 & 索宝蛋白 & 粮油加工 & 18.2 & 9.7 & 24.9 & 0.68 & 3.4 & 13.6 \\
603717 & 天域生物 & 生猪养殖 & 17.2 & 3.5 & -15.0 & -0.54 & -14.6 & 81.0 \\
002982 & 湘佳股份 & 肉鸡养殖 & 16.2 & 23.2 & 8.6 & -0.12 & -0.7 & 60.9 \\
300268 & 佳沃食品 & 其他农产品加工 & 16.0 & 4.2 & -66.1 & -0.04 & -1.4 & 35.5 \\
\end{longtable}
\endgroup


\vspace{10pt}
\begin{table}[htbp]\centering\footnotesize
\begin{tabular}{@{}lrrrrrr@{}}
\toprule
申万三级行业 & 公司数 & 总流通市值 & 平均营收 & 营收 CAGR & 亏损 & 平均 ROE \\
 & & （亿元） & 2025（亿元） & 中位数（\%） & 公司数 & 2025（\%） \\
\midrule
其他农产品加工 & 11 & 716 & 60.5 & -7.3 & 3 & 3.0 \\
畜禽饲料 & 11 & 815 & 213.3 & -5.2 & 7 & -5.5 \\
生猪养殖 & 11 & 2,962 & 275.6 & 13.5 & 5 & -5.6 \\
动物保健Ⅲ & 9 & 481 & 20.4 & 0.9 & 2 & 23.9 \\
肉鸡养殖 & 8 & 546 & 93.0 & 3.2 & 2 & 1.2 \\
种子 & 8 & 515 & 31.5 & -2.1 & 1 & 2.9 \\
粮油加工 & 7 & 373 & 387.4 & -5.9 & 1 & 4.2 \\
其他种植业 & 6 & 559 & 94.7 & 7.6 & 3 & -2.9 \\
果蔬加工 & 5 & 375 & 13.9 & 15.1 & 2 & 5.7 \\
林业Ⅲ & 4 & 244 & 5.0 & -15.6 & 3 & -4.9 \\
水产饲料 & 4 & 815 & 362.5 & 1.2 & 1 & 3.2 \\
食用菌 & 4 & 108 & 13.6 & 9.1 & 0 & 5.4 \\
水产养殖 & 4 & 119 & 16.1 & -14.8 & 2 & -45.5 \\
其他养殖 & 3 & 95 & 22.7 & 20.7 & 2 & -1.9 \\
宠物食品 & 3 & 189 & 44.8 & 18.0 & 0 & 11.5 \\
农业综合Ⅱ & 2 & 81 & 94.5 & -1.8 & 0 & 3.7 \\
海洋捕捞 & 2 & 59 & 33.4 & 5.7 & 0 & 12.2 \\
粮食种植 & 2 & 348 & 76.5 & -3.6 & 0 & 11.4 \\
\bottomrule
\end{tabular}
\caption{按申万三级行业汇总的农业上市公司数量与经营指标（2025 年年报口径；ROE 为算术平均，市值时点 2026-09-24）}
\label{tab:company-by-industry}
\end{table}


\vspace{10pt}
\begingroup\scriptsize\setlength{\tabcolsep}{1.5pt}
\begin{longtable}{@{}llp{1.9cm}rrrrl@{}}
\caption{农业上市公司债务结构与滚动偿付压力（2026 年 6 月末；按现金短债比升序；单位：亿元，现金短债比为倍）}\label{tab:company-financing}\\
\toprule
代码 & 公司简称 & 申万三级行业 & 货币资金 & 短期债务 & 有息负债 & 现金短债比 & 静态缺口 & 融资需求档 \\
\midrule
\endfirsthead
\multicolumn{9}{l}{\small（续）}\\
\toprule
代码 & 公司简称 & 申万三级行业 & 货币资金 & 短期债务 & 有息负债 & 现金短债比 & 静态缺口 & 融资需求档 \\
\midrule
\endhead
\bottomrule
\endlastfoot
002124 & 天邦食品 & 生猪养殖 & 1.55 & 35.01 & 48.40 & 0.04 & 33 & 第一档 \\
000019 & 深粮控股 & 粮油加工 & 0.87 & 18.04 & 21.43 & 0.05 & 17 & --- \\
000911 & *ST广糖 & 其他农产品加工 & 2.29 & 29.24 & 31.13 & 0.08 & 27 & 第一档 \\
002548 & 金新农 & 畜禽饲料 & 2.36 & 22.72 & 28.92 & 0.10 & 20 & 第一档 \\
300761 & 立华股份 & 肉鸡养殖 & 3.67 & 29.94 & 35.68 & 0.12 & 26 & --- \\
603668 & 天马科技 & 水产饲料 & 4.61 & 34.56 & 41.55 & 0.13 & 30 & --- \\
600962 & 国投中鲁 & 果蔬加工 & 1.57 & 9.37 & 10.30 & 0.17 & 8 & 第二档 \\
002299 & 圣农发展 & 肉鸡养殖 & 11.73 & 69.56 & 76.86 & 0.17 & 58 & --- \\
300511 & 雪榕生物 & 食用菌 & 1.76 & 10.20 & 12.52 & 0.17 & 8 & --- \\
002868 & 绿康生化 & 动物保健Ⅲ & 0.51 & 2.94 & 3.19 & 0.17 & 2 & 第二档 \\
300094 & 国联水产 & 水产养殖 & 1.98 & 10.29 & 10.84 & 0.19 & 8 & --- \\
600108 & 亚盛集团 & 其他种植业 & 5.31 & 27.04 & 46.72 & 0.20 & 22 & --- \\
600127 & 金健米业 & 粮油加工 & 1.50 & 7.28 & 7.31 & 0.21 & 6 & --- \\
003030 & 祖名股份 & 其他农产品加工 & 1.80 & 8.74 & 10.02 & 0.21 & 7 & --- \\
600467 & 好当家 & 水产养殖 & 5.64 & 23.80 & 27.09 & 0.24 & 18 & --- \\
300498 & 温氏股份 & 生猪养殖 & 36.75 & 146.20 & 286.24 & 0.25 & 109 & --- \\
301116 & 益客食品 & 肉鸡养殖 & 4.84 & 17.56 & 23.69 & 0.28 & 13 & --- \\
300268 & 佳沃食品 & 其他农产品加工 & 0.48 & 1.70 & 1.72 & 0.28 & 1 & 第二档 \\
600737 & 中粮糖业 & 其他农产品加工 & 8.58 & 30.28 & 30.71 & 0.28 & 22 & 第三档 \\
001201 & 东瑞股份 & 生猪养殖 & 4.55 & 15.69 & 24.27 & 0.29 & 11 & --- \\
002714 & 牧原股份 & 生猪养殖 & 165.47 & 568.32 & 748.02 & 0.29 & 403 & --- \\
002385 & 大北农 & 畜禽饲料 & 32.81 & 104.33 & 135.45 & 0.31 & 72 & 第一档 \\
000702 & 正虹科技 & 畜禽饲料 & 0.60 & 1.81 & 1.97 & 0.33 & 1 & --- \\
600191 & 华资实业 & 其他农产品加工 & 0.46 & 1.38 & 2.53 & 0.33 & 1 & --- \\
603477 & 巨星农牧 & 生猪养殖 & 6.94 & 20.13 & 37.45 & 0.34 & 13 & 第一档 \\
002321 & 华英农业 & 其他养殖 & 1.92 & 5.52 & 6.13 & 0.35 & 4 & 第一档 \\
002696 & 百洋股份 & 水产饲料 & 5.32 & 15.11 & 16.70 & 0.35 & 10 & 第一档 \\
002215 & 诺普信 & 其他种植业 & 12.33 & 33.99 & 58.48 & 0.36 & 22 & 第三档 \\
000998 & 隆平高科 & 种子 & 24.25 & 64.83 & 107.27 & 0.37 & 41 & --- \\
000876 & 新希望 & 畜禽饲料 & 79.63 & 205.48 & 579.09 & 0.39 & 126 & 第一档 \\
002852 & 道道全 & 粮油加工 & 4.92 & 12.69 & 12.69 & 0.39 & 8 & --- \\
002069 & 獐子岛 & 水产养殖 & 5.45 & 14.05 & 19.24 & 0.39 & 9 & --- \\
300087 & ST荃银 & 种子 & 7.67 & 19.06 & 22.07 & 0.40 & 11 & 第二档 \\
002556 & 辉隆股份 & --- & 11.20 & 27.68 & 31.42 & 0.40 & 16 & --- \\
000735 & 罗牛山 & 生猪养殖 & 4.11 & 9.97 & 19.18 & 0.41 & 6 & --- \\
002234 & 民和股份 & 肉鸡养殖 & 5.14 & 12.38 & 13.65 & 0.42 & 7 & --- \\
300119 & 瑞普生物 & 动物保健Ⅲ & 5.32 & 12.68 & 16.16 & 0.42 & 7 & 第二档 \\
002481 & 双塔食品 & 其他农产品加工 & 3.84 & 8.95 & 12.29 & 0.43 & 5 & --- \\
300313 & 天山生物 & 其他养殖 & 0.26 & 0.58 & 0.93 & 0.45 & 0 & --- \\
000048 & 京基智农 & 生猪养殖 & 6.08 & 13.29 & 19.75 & 0.46 & 7 & --- \\
603363 & 傲农生物 & 畜禽饲料 & 4.75 & 9.72 & 23.80 & 0.49 & 5 & --- \\
603717 & 天域生物 & 生猪养殖 & 1.35 & 2.75 & 10.49 & 0.49 & 1 & 第一档 \\
600975 & 新五丰 & 生猪养殖 & 6.08 & 12.10 & 73.49 & 0.50 & 6 & 第一档 \\
002688 & 金河生物 & 动物保健Ⅲ & 7.00 & 13.82 & 23.64 & 0.51 & 7 & 第三档 \\
000798 & 中水渔业 & 海洋捕捞 & 8.75 & 16.75 & 27.86 & 0.52 & 8 & --- \\
300138 & 晨光生物 & 其他农产品加工 & 28.02 & 52.71 & 53.17 & 0.53 & 25 & 第三档 \\
603609 & 禾丰股份 & 畜禽饲料 & 12.18 & 22.89 & 66.24 & 0.53 & 11 & --- \\
300999 & 金龙鱼 & 粮油加工 & 485.84 & 887.12 & 927.22 & 0.55 & 401 & --- \\
601118 & 海南橡胶 & 其他种植业 & 51.31 & 93.23 & 167.80 & 0.55 & 42 & 第一档 \\
603151 & 邦基科技 & 畜禽饲料 & 7.22 & 13.06 & 15.02 & 0.55 & 6 & 第二档 \\
600540 & 新赛股份 & 其他种植业 & 8.39 & 14.35 & 15.08 & 0.58 & 6 & --- \\
000972 & 中基健康 & 果蔬加工 & 3.01 & 5.01 & 5.54 & 0.60 & 2 & 第一档 \\
000930 & 中粮科技 & 其他农产品加工 & 21.29 & 34.61 & 49.14 & 0.62 & 13 & --- \\
605296 & 神农集团 & 生猪养殖 & 4.50 & 6.59 & 13.62 & 0.68 & 2 & --- \\
300021 & 大禹节水 & --- & 12.85 & 18.48 & 29.12 & 0.70 & 6 & 第一档 \\
600359 & 新农开发 & 其他种植业 & 3.42 & 4.87 & 6.52 & 0.70 & 1 & --- \\
688156 & 路德科技 & 畜禽饲料 & 2.43 & 3.38 & 5.50 & 0.72 & 1 & --- \\
002679 & 福建金森 & 林业Ⅲ & 1.34 & 1.75 & 11.56 & 0.77 & 0 & --- \\
000592 & 平潭发展 & 林业Ⅲ & 4.17 & 5.27 & 5.30 & 0.79 & 1 & --- \\
002100 & 天康生物 & 畜禽饲料 & 25.35 & 30.51 & 65.33 & 0.83 & 5 & --- \\
002311 & 海大集团 & 水产饲料 & 34.25 & 40.97 & 64.39 & 0.84 & 7 & 第三档 \\
002567 & 唐人神 & 畜禽饲料 & 21.10 & 24.84 & 61.18 & 0.85 & 4 & 第一档 \\
002891 & 中宠股份 & 宠物食品 & 13.68 & 16.07 & 18.56 & 0.85 & 2 & --- \\
002458 & 益生股份 & 肉鸡养殖 & 14.79 & 16.42 & 18.88 & 0.90 & 2 & 第三档 \\
002982 & 湘佳股份 & 肉鸡养殖 & 3.94 & 4.02 & 19.28 & 0.98 & 0 & --- \\
600257 & 大湖股份 & 水产养殖 & 2.61 & 2.50 & 6.87 & 1.04 & -0 & --- \\
301498 & 乖宝宠物 & 宠物食品 & 9.85 & 9.29 & 11.94 & 1.06 & -1 & --- \\
001366 & 播恩集团 & 畜禽饲料 & 2.60 & 2.37 & 3.68 & 1.10 & -0 & --- \\
300189 & 神农种业 & 种子 & 0.99 & 0.88 & 3.05 & 1.11 & -0 & --- \\
600251 & 冠农股份 & 果蔬加工 & 13.50 & 11.36 & 11.87 & 1.19 & -2 & 第三档 \\
001313 & 粤海饲料 & 水产饲料 & 6.78 & 5.52 & 8.40 & 1.23 & -1 & --- \\
688098 & 申联生物 & 动物保健Ⅲ & 0.65 & 0.53 & 0.53 & 1.23 & -0 & --- \\
002772 & 众兴菌业 & 食用菌 & 16.69 & 12.38 & 35.65 & 1.35 & -4 & 第三档 \\
002746 & 仙坛股份 & 肉鸡养殖 & 36.20 & 25.39 & 25.92 & 1.43 & -11 & --- \\
300967 & 晓鸣股份 & 肉鸡养殖 & 1.43 & 0.96 & 6.43 & 1.49 & -0 & --- \\
000505 & 京粮控股 & 粮油加工 & 16.90 & 10.88 & 17.91 & 1.55 & -6 & --- \\
300871 & 回盛生物 & 动物保健Ⅲ & 3.59 & 2.02 & 4.64 & 1.77 & -2 & 第三档 \\
001338 & 永顺泰 & 其他农产品加工 & 14.14 & 7.38 & 7.41 & 1.92 & -7 & --- \\
600195 & 中牧股份 & 动物保健Ⅲ & 16.65 & 8.44 & 11.12 & 1.97 & -8 & --- \\
600097 & 开创国际 & 海洋捕捞 & 3.90 & 1.76 & 4.79 & 2.22 & -2 & --- \\
002286 & 保龄宝 & 其他农产品加工 & 3.13 & 1.37 & 1.77 & 2.29 & -2 & --- \\
603336 & 宏辉果蔬 & 其他种植业 & 4.04 & 1.71 & 7.43 & 2.37 & -2 & --- \\
002157 & 正邦科技 & 生猪养殖 & 16.97 & 6.97 & 19.10 & 2.43 & -10 & --- \\
600313 & 农发种业 & 种子 & 8.73 & 3.58 & 4.77 & 2.44 & -5 & --- \\
601952 & 苏垦农发 & 粮食种植 & 14.56 & 4.28 & 45.25 & 3.40 & -10 & --- \\
603231 & 索宝蛋白 & 粮油加工 & 2.84 & 0.65 & 0.67 & 4.39 & -2 & --- \\
000663 & 永安林业 & 林业Ⅲ & 1.68 & 0.32 & 0.55 & 5.28 & -1 & 第二档 \\
000713 & 国投丰乐 & 种子 & 15.50 & 2.38 & 3.72 & 6.52 & -13 & --- \\
603566 & 普莱柯 & 动物保健Ⅲ & 0.86 & 0.07 & 0.26 & 11.86 & -1 & --- \\
300970 & 华绿生物 & 食用菌 & 5.89 & 0.31 & 7.00 & 19.04 & -6 & --- \\
600354 & 敦煌种业 & 种子 & 8.05 & 0.33 & 0.65 & 24.20 & -8 & 第三档 \\
600883 & 博闻科技 & 食用菌 & 1.43 & 0.05 & 0.12 & 30.40 & -1 & --- \\
300673 & 佩蒂股份 & 宠物食品 & 4.95 & 0.14 & 8.21 & 34.24 & -5 & --- \\
002086 & 东方海洋 & 其他农产品加工 & 2.43 & 0.06 & 0.49 & 37.76 & -2 & --- \\
600965 & *ST福成 & 其他养殖 & 2.03 & 0.05 & 0.98 & 43.48 & -2 & --- \\
600265 & *ST景谷 & 林业Ⅲ & 0.92 & 0.02 & 0.04 & 55.11 & -1 & 第二档 \\
300175 & 朗源股份 & 果蔬加工 & 0.84 & 0.01 & 0.28 & 61.19 & -1 & --- \\
688526 & 科前生物 & 动物保健Ⅲ & 1.70 & 0.02 & 0.07 & 96.61 & -2 & --- \\
002041 & 登海种业 & 种子 & 21.43 & 0.19 & 0.26 & 112.76 & -21 & --- \\
600598 & 北大荒 & 粮食种植 & 19.27 & 0.03 & 0.10 & 566.76 & -19 & --- \\
600201 & 生物股份 & 动物保健Ⅲ & 16.09 & 0.01 & 1.01 & 1,426.53 & -16 & --- \\
605198 & 安德利 & 果蔬加工 & 7.39 & --- & --- & --- & --- & --- \\
600371 & 万向德农 & 种子 & 2.26 & --- & 0.00 & --- & --- & --- \\
603182 & 嘉华股份 & 粮油加工 & 1.39 & --- & 0.14 & --- & --- & --- \\
\end{longtable}
\endgroup


\vspace{10pt}
\noindent\textbf{使用说明}：表中“同比”为 2026 年上半年营业收入相对 2025 年上半年的增速；
“ROE”为半年报披露的净资产收益率，\hl{未年化}（如需与年报口径比较，可近似乘以 2），
净资产为负的公司该指标记为“—”；
“—”表示该数据在所获取的公开接口中缺失。
全部数据由 \texttt{scripts/make\_company.py} 生成（正文统计数据由
\texttt{scripts/company\_stats.py} 汇总），
可通过 \texttt{data/clean/company\_master.csv} 复核任意单元格。
